\documentclass[11pt,a4paper]{article}
\usepackage[T1]{fontenc}
\usepackage[utf8]{inputenc}
\usepackage[margin=2.5cm]{geometry}
\usepackage{xcolor}
\usepackage{parskip}
\usepackage[hidelinks]{hyperref}
\newcommand{\tbc}[1]{\textcolor{red}{[TO BE COMPLETED: #1]}}
\pagestyle{empty}

\begin{document}

\begin{flushright}
Bartosz Lewandowski\\
Department of Clinical Engineering\\
Academy of Silesia in Katowice, Katowice, Poland\\
\href{mailto:bartosz.lewandowski@akademiaslaska.pl}{bartosz.lewandowski@akademiaslaska.pl}\\[6pt]
\today
\end{flushright}

Editorial Office\\
\textit{Foods} (MDPI)

Dear Editor,

On behalf of all co-authors, I submit the manuscript entitled
\textbf{``Reproducibility of Honey Taste and Odor Typicality Assessment: A Blind
Consumer Panel Versus the Verdicts of Four Laboratories''} for consideration
as a research article in \textit{Foods}, section  \textit{Food Quality and Safety}.

Under Polish law, honey must have an odor and taste ``typical of the given
variety'', and an organoleptic verdict of atypical taste can by itself lead to
the rejection of a batch. Neither EU nor Polish law defines a reference for
typicality, and the reproducibility of such verdicts is rarely tested. We
assessed the same seven honey batches in a blind consumer panel
(138~respondents, 966~ratings) and, in parallel, by four commercial and
institutional laboratories, and we quantified the agreement at three levels:
the panel as a rater, pairs of raters (laboratory--laboratory,
consumer--consumer, consumer--laboratory), and the ability of a single rating
to reproduce the laboratory verdict.

The main findings are:
\begin{itemize}
\item agreement among the four laboratories was fair on the binary scale
($\kappa = 0.30$) and slight on the three-level scale ($\kappa = 0.13$);
\item a consumer agreed with a laboratory to the same degree as with another
consumer ($\kappa = 0.10$ and $0.08$), and a single consumer rating reproduced
the laboratory majority verdict with $\mathrm{AUC} = 0.57$;
\item the panel and the laboratories agreed on the clearly typical sample, but
not on the borderline ones.
\end{itemize}
These results bear directly on the use of organoleptic typicality as a
stand-alone criterion in the official and commercial control of honey quality,
a topic within the scope of \textit{Foods}.

We confirm that this manuscript has not been published and is not under consideration 
for publication elsewhere, and that all authors have approved the manuscript 
and agree with its submission. The authors declare no conflicts of interest.



Yours sincerely,

Bartosz Lewandowski\\
(corresponding author, on behalf of all authors)

\end{document}
